% mode: plain
% plain.tex has \outer macros, so here the flag is on from the format: every control sequence is looked up, and the outer ones stop a loop.
\tracingonline=1 \errorcontextlines=100 \tracingmacros=1
\def\a#1{[#1]}
\def\c#1.{[#1]}
\def\s#1{\iffalse#1\fi}
\def\name#1{\expandafter\noexpand\csname#1\endcsname}
\def\mk#1{\expandafter\def\csname in#1\expandafter\endcsname\expandafter{\expandafter\a\expandafter{\csname#1\endcsname}}}
\mk{newcount}
\show\innewcount
\message{^^Jan outer macro in a group argument, from a list:}
\innewcount
\expandafter\def\expandafter\ib\expandafter{\expandafter\c\expandafter x\csname bye\endcsname.}
\message{^^Jthe same in a delimited argument:}
\ib
\expandafter\def\expandafter\ic\expandafter{\expandafter\iffalse\csname newif\endcsname\fi}
\message{^^Jand skipped:}
\ic
\fi
\expandafter\def\expandafter\id\expandafter{\csname newdimen\endcsname}
\message{^^Jread in place:}
\expandafter\a\expandafter{\id}
\expandafter\s\expandafter{\id}
\fi
\message{^^Jthe usual macros, with the flag on:}
\def\ok{\a{x}\c y.\s{z}\edef\e{\a{x}}\show\e}
\ok
\newcount\mycount
\newif\ifmine
\minetrue
\ifmine\message{mine}\else\message{not mine}\fi
\minefalse
\ifmine\message{mine}\else\message{not mine}\fi
\def\fam#1{\setbox0\hbox{#1}\showbox0}
\fam{\vbox{\halign{\strut#\cr a\cr b\cr}}}
\fam{\vbox{\settabs 2\columns \+ a&b\cr}}
\fam{$\matrix{a&b\cr c&d\cr}$}
\beginsection A section\par
\proclaim Theorem. Something.\par
\message{^^Jand at the end of the job:}
\bye
